\section{Exact probability normalization and projection identities}
\label{or:aor-probability-normalization}

For a finite carrier \(A\), every probability law in the model is an exact
rational table
\[
  \Delta_{\mathbb Q}(A):=
  \left\{p:A\to\mathbb Q:
  p(a)\geq0\ \text{for every }a,
  \sum_{a\in A}p(a)=1\right\}.
\]
Beliefs, belief-kernel rows, and candidate-generation rows use this convention.
At fixed project \(q\), belief \(b\), and closure \(C\), verification assigns
\begin{align}
 \Gamma(q,b,C)(s)
   &=G(q,b,C)(s)\nu(q,b,C,s),\qquad s\in S,
   \label{eq:or-admission-some}\\
 \Gamma(q,b,C)(\mathsf{none})
   &=G(q,b,C)(\mathsf{none})
     +\sum_{s\in S}G(q,b,C)(s)
       \bigl(1-\nu(q,b,C,s)\bigr).
   \label{eq:or-admission-none}
\end{align}
Every term is nonnegative. Summing
\eqref{eq:or-admission-some}--\eqref{eq:or-admission-none} cancels the
verification split and recovers the unit mass of \(G\), so \(\Gamma\) is
normalized.

Conditional on initial belief \(b\), productive state \(K=(F,C)\), and
project \(q\) of duration \(d_q\), let
\(\Lambda_q(\mathbf b,o\mid b,K)\) be the exact joint law of the belief path
\(\mathbf b=(b_0,\ldots,b_{d_q})\) and admitted outcome \(o\). Its specified
marginals are
\begin{align}
 \sum_o\Lambda_q(\mathbf b,o\mid b,K)
   &=\mathbf1\{b_0=b\}
     \prod_{t=0}^{d_q-1}P_B(b_t,b_{t+1}),
   \label{eq:or-completion-path-marginal}\\
 \sum_{\mathbf b}\Lambda_q(\mathbf b,o\mid b,K)
   &=\Gamma(q,b,C)(o).
   \label{eq:or-completion-outcome-marginal}
\end{align}
These equations do not impose independence between the belief path and the
admitted outcome.

For a raw source library \(L\), deterministic update gives
\[
 \mathcal Q_q^{\mathrm{raw}}(b',L'\mid b,L)
 :=\sum_{\substack{\mathbf b,o:\ b_{d_q}=b'\\L\oplus o=L'}}
   \Lambda_q(\mathbf b,o\mid b,K_L).
\]
For realizable \(K=(F,C)\), the two compressed pushforwards are
\begin{align}
 \overline T_q(K'\mid b,K)
   &:=\sum_{o:\operatorname{addK}(K,o)=K'}
       \Gamma(q,b,C)(o),
   \label{eq:or-compressed-outcome-law}\\
 \overline{\mathcal Q}_q(b',K'\mid b,K)
   &:=\sum_{\substack{\mathbf b,o:\ b_{d_q}=b'\\
       \operatorname{addK}(K,o)=K'}}
       \Lambda_q(\mathbf b,o\mid b,K).
   \label{eq:or-compressed-joint-law}
\end{align}
Finite deterministic pushforward preserves nonnegativity and total mass.
The local update identity
\(K_{L\oplus o}=\operatorname{addK}(K_L,o)\) also makes every compressed
successor realizable. These identities supply the normalization clause of the
raw-to-compressed projection theorem; they do not supply the factorization,
positive-duration induction, or contraction arguments by themselves.

\subsection{Complete projection proof}

The journal article assumes a finite hidden state, finite filtered belief
carrier, finite strategy catalog, finite project catalog, positive integer
project durations, exact rational transition and completion laws, and a
mandatory inactive strategy. Every operating reward, project availability
indicator, project cost, duration, generation law, verification law, and
joint terminal belief--admission law factors through the current belief and
the productive library state \(K_L=(F_L,C_L)\). Raw insertion and compressed
insertion obey
\begin{equation}
 K_{L\oplus o}=\operatorname{addK}(K_L,o).
 \label{eq:or-local-update}
\end{equation}
Stationary claims additionally assume a common discount factor
$0<\beta<1$; finite-horizon claims do not.

\begin{theorem}[Raw-to-compressed controlled semi-Markov projection]
\label{thm:or-raw-compressed-projection}
Under the preceding assumptions, $K_L=(F_L,C_L)$, augmented during an
unfinished project by project identity and remaining duration, is a
controlled semi-Markov state for the raw library process. Raw and compressed
finite-horizon values agree at every raw state, and the complete sets of
maximizing actions agree under the natural action identification. Under the
stationary discount clause, the raw and compressed Bellman operators have
unique fixed points with the same values and maximizing stationary-action
correspondence.
\end{theorem}

\begin{proof}
Fix a raw library $L$, belief $b$, and feasible action. For Continue,
the operating reward and belief transition depend on $L$ only through
$F_L$, hence through $K_L$. For a project $q$, the factorization
assumption makes feasibility, cost, duration, the incumbent operating-reward
block, and the joint completion law functions of $(b,K_L,q)$. Equations
\eqref{eq:or-admission-some}--\eqref{eq:or-completion-outcome-marginal}
show that this law is normalized without imposing independence between the
belief path and admission. Pushing it through the deterministic raw update
and using \eqref{eq:or-local-update} gives exactly the compressed transition
law \eqref{eq:or-compressed-joint-law}. Its support consists of realizable
compressed states. Thus two raw libraries in the same $K$-fiber have the
same feasible actions, current reward blocks, durations, and conditional laws
of the next compressed decision state.

For a finite calendar horizon $h$, prove equality of every corresponding
action value by strong induction on $h$. The claim is immediate at horizon
zero. Continue invokes the induction hypothesis at $h-1$; a feasible
project invokes it after its positive duration $d_q$, at $h-d_q<h$.
Current rewards and transition laws already agree by the preceding paragraph.
Each corresponding action value therefore agrees, as do the finite maxima and
the complete maximizing-action sets. Finiteness makes every maximum attained.

For the stationary clause, cast the exact rational tables into the reals and
use the sup norm on the finite state carrier. Conditional expectation is
nonexpansive. Continue discounts continuation by $\beta$, and a project
discounts it by $\beta^{d_q}\leq\beta$. Maximization over the common finite
action set preserves this modulus, so both Bellman operators are
$\beta$-contractions. Bellman intertwining, established by the action-value
identity above, maps the compressed fixed point to a raw fixed point. Banach
uniqueness gives value equality; the same pointwise action-value identity gives
the complete maximizing stationary-action correspondence. No closure
detectability assumption enters this forward projection.
\end{proof}
